% !PRINTED 56
\weberpage{56}

11. \gesperrt{Numerator and denominator of this representation are uniquely determined by $\varphi(u)$ itself, apart from unit factors.}

The functions $G(u)$, $G_1(u)$ can be decomposed into prime factors, i.e. into factors which become zero at only one point. If one pursues this path in the case of the doubly periodic functions, one arrives at the \gesperrt{Weierstrass} $\sigma$-functions. We shall, however, take another path, on which we shall not meet the $\sigma$-functions until a later place.

\begin{center}
\textbf{\S\,16. Periodicity.}
\end{center}

A function $\varphi(u)$ of $u$ which has the property that, for a constant $\omega$ and for every value of $u$,
\weq{1}{\varphi(u+\omega) = \varphi(u)}
holds, is called \gesperrt{periodic}, and $\omega$ is called the \gesperrt{period}.

The subject of our investigation is functions $\varphi(u)$ with two periods $\omega_1, \omega_2$, of which we shall presuppose once for all that their ratio $\omega_2:\omega_1$ is not real and that the imaginary part of this ratio is \gesperrt{positive}. In the latter assumption there lies only a stipulation concerning the notation $\omega_1, \omega_2$. For if $\omega_2:\omega_1$ has a negative imaginary part, then that of $\omega_1:\omega_2$ is positive.

Every combination $m_1\omega_1 + m_2\omega_2$ is likewise a period, if $m_1, m_2$ are integers.

We take an arbitrary point $u_0$ in the plane $u$ and denote the four points
\[
u_0, \quad u_0 + \omega_1, \quad u_0 + \omega_1 + \omega_2, \quad u_0 + \omega_2;
\]
if we join these four points in order by straight lines, returning from the last to the first, we obtain a parallelogram, which is called the \gesperrt{period parallelogram}. The rectilinear sides of the period parallelogram may also be replaced by curvilinear paths, provided only that the opposite ones come into coincidence by parallel displacement and that the individual paths form no loops.

If $\omega_1 = \alpha_1 + \beta_1 i$, $\omega_2 = \alpha_2 + \beta_3 i$ and $\alpha_1, \alpha_2, \beta_1, \beta_3$ are real, % [sic] German source prints beta_3 twice; mathematically expected beta_2.
then by hypothesis $(\alpha_1\beta_2 - \alpha_2\beta_1)$ is positive, and the geometrical meaning of this quantity is (by known theorems
%CONT>>
